%% Springer Nature journal article template (sn-jnl)
%% Target journal: Nature Genetics (Article / Technical Report)
%% Nature Portfolio reference style
%%
%% PAPER #2 — Graph-native epistasis (placeholder skeleton)
%% Companion to paper #1 (paper/finemapping_v1/) on relational fine-mapping.

\documentclass[pdflatex,sn-nature]{sn-jnl}

%%% Standard Packages
\usepackage{graphicx}
\usepackage{multirow}
\usepackage{amsmath,amssymb,amsfonts}
\usepackage{booktabs}
\usepackage{xcolor}
\usepackage{textcomp}
\usepackage{hyperref}
\usepackage{siunitx}
\usepackage{array}
\usepackage{longtable}
\usepackage{enumitem}
\usepackage{lineno}
\linenumbers

\graphicspath{{figures/}}

\raggedbottom
\unnumbered

\begin{document}

\title[Graph-native epistasis]{Graph-native epistasis detection rescues anti-conservative bias in biobank-scale GWAS}

%%=============================================================%%
%% Author information
%%=============================================================%%

\author[1]{\fnm{Ehsan} \sur{Estaji}}\email{ehsan.estaji@umu.se}
\equalcont{These authors contributed equally to this work.}

\author[1]{\fnm{Shi-Wei} \sur{Zhao}}\email{shi-wei.zhao@umu.se}
\equalcont{These authors contributed equally to this work.}

\author[1]{\fnm{Zhao-Yang} \sur{Chen}}\email{zhao-yang.chen@umu.se}

\author[2]{\fnm{Shuai} \sur{Nie}}\email{nieshuai@gdaas.cn}

\author*[1]{\fnm{Jian-Feng} \sur{Mao}}\email{jianfeng.mao@umu.se}

\affil[1]{\orgdiv{Ume\aa{} Plant Science Centre, Department of Plant Physiology}, \orgname{Ume\aa{} University}, \orgaddress{\postcode{SE-90187}, \city{Ume\aa{}}, \country{Sweden}}}

\affil[2]{\orgdiv{Rice Research Institute, Key Laboratory of Genetics and Breeding of High Quality Rice in Southern China (Co-construction by Ministry and Province), Ministry of Agriculture and Rural Affairs, Guangdong Key Laboratory of New Technology in Rice Breeding}, \orgname{Guangdong Academy of Agricultural Sciences}, \orgaddress{\postcode{510640}, \city{Guangzhou}, \country{China}}}

% PLACEHOLDER abstract — Yelmen-reframing pitch.  Final empirical numbers will
% replace [PLACEHOLDER] tags as Phase C/D experiments produce results.

\abstract{
Classical genome-wide association studies --- including state-of-the-art
linear mixed models such as REGENIE --- are anti-conservative when the
true genetic architecture involves epistasis but the analytical model
omits it. Yelmen \emph{et~al.}\ (2026) showed that the bias grows with
sample size, so biobank-scale GWAS is more vulnerable, not less.
We present GraphGWAS-Epistasis, a graph-native epistasis-detection
toolkit that exploits a typed multi-omics knowledge graph (genes,
tissue-specific eQTLs, pathways, protein--protein interactions) to
build a biology-typed interaction subspace. Four detection strategies
operate on this substrate: M1 LD-pruned co-occurrence with explicit
interaction tests; M2 motif-filtered pairs (same-gene, PPI-partner,
same-pathway); M3 differential subgraph analysis between cases and
controls; and M4 dark-matter / synthetic-incompatibility detection.
On simulated phenotypes derived from [PLACEHOLDER --- 1KG chr21+22 /
Estonian Biobank if access granted], feeding the detected interactions
back into the LMM as covariates collapses the spurious-significance
regime from [PLACEHOLDER]\,\% to [PLACEHOLDER]\,\%. The same graph-typed
interaction subspace yields a closed-form $\rho_{\max}$ diagnostic that
flags bias-prone target SNPs before testing. Three higher-order methods
(M5 random walk, M6 GNN-attention, M7 persistent homology) extend the
framework to $k\geq 3$ interactions that pairwise tests cannot
formulate. The toolkit is released as part of the GraphGWAS Python
package; all benchmarks reproduce from a single command.
}


\keywords{epistasis, gene-gene interaction, GWAS bias, linear mixed model, biobank, multi-omics, graph database, REGENIE}

\maketitle

% Main text
% PLACEHOLDER introduction — to expand as Phase C/D results land.

\section{Introduction}\label{sec:intro}

[PLACEHOLDER paragraph 1 --- the bias problem.] Classical GWAS, even
with state-of-the-art linear mixed models, omits gene--gene interaction
terms. Yelmen \emph{et~al.}\ \cite{yelmen2026} recently characterised
the resulting bias mathematically: the null $t$-statistic acquires a
mean shift $\mu = \rho\sqrt{\lambda n}/\sqrt{1-\lambda\rho^{2}}$ that
grows with sample size $n$, where $\rho$ is the correlation between
the target SNP and the realised interaction signal and $\lambda$ is
the variance fraction explained by interactions. At biobank scale
($n\geq 100{,}000$) this shift produces 1{,}000s of spurious
genome-wide significant hits even under their strict no-path null
where the target SNP cannot itself participate in the interaction.

[PLACEHOLDER paragraph 2 --- the detection-tool landscape.] Existing
epistasis-detection tools (BOOST \cite{wan2010boost},
MAPIT \cite{crawford2017}, MDR-family approaches) treat the search space
as either an exhaustive pairwise grid (computationally infeasible at
biobank scale) or an unconstrained statistical fishing exercise (low
power after multiple-testing correction). Neither leverages the
biology-typed structure that genome-wide functional resources
(GENCODE, GTEx, STRING, ENCODE, Reactome) make available.

[PLACEHOLDER paragraph 3 --- our contribution.] We present
GraphGWAS-Epistasis, a graph-native toolkit that uses the typed
multi-omics knowledge graph from \cite{estaji2026graphgwas} to
restrict epistasis search to biologically motivated motifs. Four
pairwise detection strategies (M1--M4) operate on this graph-typed
interaction subspace; three higher-order extensions (M5--M7) detect
$k\geq 3$ interactions. The framework also provides a closed-form
$\rho_{\max}$ diagnostic that flags bias-prone target SNPs before
testing.

[PLACEHOLDER paragraph 4 --- empirical headlines, to be filled.]

% PLACEHOLDER results — planned-experiment headers; numbers fill in as Phase C/D progresses.

\section{Results}\label{sec:results}

\subsection{Graph-typed $\rho_{\max}$ collapses bias under realistic interaction architectures}\label{sec:rho_max}

[PLACEHOLDER --- replicates Yelmen \emph{et~al.}\ Fig~1 on 1KG~chr21+22
under random $Z$ (mean $\rho_{\max}\approx 0.032$, max~$\approx 0.04$
across chromosomes; max~$0.85$ within-chromosome). Then extends to
motif-typed $Z$: shows that for genuine causal target SNPs, motif-typed
$\rho_{\max}$ is suppressed; for null target SNPs, it stays comparable
to random. Empirical claim: graph-typed interaction modelling reduces
the bias selectively where it matters.]

\subsection{Conservativeness rescue on REGENIE-LMM-detected loci}\label{sec:rescue}

[PLACEHOLDER --- run REGENIE on simulated phenotypes (1KG chr21+22 +
realistic $\lambda$ values); count spuriously significant hits per
\cite{yelmen2026}'s criterion. Then add GraphGWAS M2-detected
interactions as covariates; recount. Headline number
$X{\to}Y$~percent reduction in spurious hits.]

\subsection{Power vs BOOST, MAPIT, MDR}\label{sec:power_baselines}

[PLACEHOLDER --- 100 replicates per scenario S1--S5; rank-1 detection
rate, mean rank, runtime per locus. Expectation: M2 motif filtering
trades raw recall for prior plausibility, beating exhaustive baselines
on FDR at low $\beta_{I}$.]

\subsection{Yeast real-data validation}\label{sec:yeast_validation}

[PLACEHOLDER --- 1011 Yeast Genomes; literature-validated epistatic
pairs (BCY1--TPK1, HSP104 modifiers, RAS/PKA crosstalk).]

\subsection{Higher-order interactions ($k\geq 3$)}\label{sec:higher_order}

[PLACEHOLDER --- M5/M6/M7 demonstration; the structurally unique claim
that pairwise methods cannot reach.]

\subsection{Cross-species replication}\label{sec:cross_species}

[PLACEHOLDER --- \emph{Arabidopsis} 1001~Genomes, 3K~Rice~Genomes,
human~Pan-UKB.]

% PLACEHOLDER discussion — to fill once empirical results are in.

\section{Discussion}\label{sec:discussion}

[PLACEHOLDER paragraph 1 --- recap of the bias problem and our
contribution.]

[PLACEHOLDER paragraph 2 --- limitations: $\rho$ is unobservable
in real GWAS, so the $\rho_{\max}$ upper bound is what we use;
our M2 motif filtering is biased toward known biology, so novel
interactions outside known pathways/PPIs are under-detected.]

[PLACEHOLDER paragraph 3 --- biobank applications: with EstBB / UKB
access this approach can flag known-spurious-by-bias hits in published
loci.]

[PLACEHOLDER paragraph 4 --- forward-looking: the bias diagnostic
formalism extends naturally to gene--environment interactions
(\cite{herrera2024gxe}) and higher-order ($k\geq 4$) interactions.]


% Online Methods
\section{Online Methods}\label{sec:methods}
% PLACEHOLDER methods — section headers + brief intent for each subsection.
% Implementation status (M1-M4 done; M5-M7 designed) tracked in
% docs/EPISTASIS_V2_DESIGN.md.

\subsection{Graph-native epistasis substrate}\label{sec:methods_substrate}

[PLACEHOLDER --- recap from paper~\#1 \cite{estaji2026graphgwas}: typed
multi-omics knowledge graph (Variant--Gene--Pathway--Tissue), schema,
node and edge counts on 1000 Genomes / Pan-UK Biobank.]

\subsection{LD-pruned co-occurrence with interaction test (M1)}\label{sec:methods_m1}

[PLACEHOLDER --- LD-aware maximum independent set on the LD graph at
threshold $\tau$; on survivors, fit $y \sim G_{1} + G_{2} + G_{1}G_{2}$
with logistic / linear / Firth correction; BH-FDR. Theorem~1
(search-space reduction $\sim 1/k(\tau)^{2}$) recapped from
\cite{estaji2026graphgwas} Supplementary Note~S1.]

\subsection{Motif-filtered pairwise testing (M2)}\label{sec:methods_m2}

[PLACEHOLDER --- candidate enumeration via Cypher motifs P1
(same-gene), P2 (PPI partner), P3 (same-pathway). Same interaction
regression as M1 on the candidate set.]

\subsection{Differential subgraph analysis (M3)}\label{sec:methods_m3}

[PLACEHOLDER --- co-occurrence sub-graphs for cases and controls;
log-fold-ratio scoring; Fisher's exact for significance.]

\subsection{Synthetic incompatibility (M4)}\label{sec:methods_m4}

[PLACEHOLDER --- depleted co-occurrence detection; Poisson tail
$p$-values; high-MAF restriction for power.]

\subsection{Higher-order interactions (M5--M7)}\label{sec:methods_m5_m7}

[PLACEHOLDER --- random walk with restart on the bipartite
variant--gene graph (M5); GNN attention-based extraction (M6);
persistent homology / 1-cycles in case-only co-occurrence
simplicial complex (M7). Implementation status table and
forthcoming-work caveats.]

\subsection{Bias diagnostic: graph-typed $\rho_{\max}$ and $R(x)$}\label{sec:methods_bias}

Following Yelmen \emph{et~al.}\ \cite{yelmen2026}, we define the
maximum correlation between a target SNP $g$ and the column space
of an interaction-feature matrix $Z$ as
\begin{equation}
\rho_{\max}(g, Z) \;=\; \frac{\lVert P_{H} g \rVert}{\lVert g \rVert},
\qquad P_{H} = Z(Z^{\top}Z)^{+}Z^{\top}.
\end{equation}
The expected $t$-statistic mean shift under the omitted-interaction
null is
\begin{equation}
\mu \;=\; \frac{\rho\sqrt{\lambda n}}{\sqrt{1-\lambda\rho^{2}}}
\end{equation}
and the conservativeness ratio
\begin{equation}
R(x) \;=\; \frac{p_{\text{true}}(x)}{p_{\text{nominal}}(x)}
\end{equation}
quantifies how anti-conservative ($R{>}1$) or conservative ($R{<}1$)
the nominal $p$-value is. These three quantities are computed by the
\texttt{graphgwas.bias} module on the same preprocessed (whitened,
covariate-residualised) feature space used by the LMM. Our
contribution beyond \cite{yelmen2026} is the substitution of a
graph-typed $Z$ (built from M2 motif pairs) for the random $Z$ they
analysed.

\subsection{Strict no-path null framework}\label{sec:methods_null}

[PLACEHOLDER --- adapted from \cite{yelmen2026} Eq.~12: target SNP
forbidden from $Z$; phenotype simulated as $y = u + \varepsilon$ with
$u \in \mathrm{col}(Z)$. Used for paper-\#2 calibration runs.]

\subsection{Cross-method validation and benchmarks}\label{sec:methods_benchmarks}

[PLACEHOLDER --- REGENIE \cite{mbatchou2021} as the LMM baseline;
BOOST, MAPIT, MDR as epistasis-specific baselines. 100~replicates
per scenario across S1 (pure interaction), S2 (marginal+interaction),
S3 (weak $\beta_{I}$), S4 (multi-pair), S5 ($k=3$). Datasets:
1KG~chr21+22 simulations; Estonian Biobank if access granted.]


\backmatter

\bmhead{Data Availability}

[PLACEHOLDER --- to be filled as Phase C results land. Likely sources:
1000~Genomes Project chr21+chr22 high-coverage NYGC release; Estonian
Biobank (subject to application); 1011 Yeast Genomes; \emph{Arabidopsis}
1001 Genomes; 3000 Rice Genomes Project; Pan-UK Biobank summary
statistics. Pre-built graph caches and benchmark JSONs to be deposited
on Zenodo.]

\bmhead{Code Availability}

GraphGWAS source code is available at
\url{https://github.com/jfmao/GraphGWAS} under the MIT licence. The
release accompanying this paper provides the Python package and
command-line interface, the epistasis-detection toolkit (M1--M7),
the bias-diagnostic module (\texttt{graphgwas.bias}), the benchmark
scripts that produced every figure and table in the paper, and a
reproducibility guide that regenerates every result from scratch.

\bmhead{Acknowledgements}

[PLACEHOLDER --- mirror finemapping\_v1 acknowledgements; add Estonian
Biobank Research Team if EstBB application is granted.]

\section*{Declarations}

\begin{itemize}
\item \textbf{Funding:} [PLACEHOLDER]
\item \textbf{Competing interests:} The authors declare no competing interests.
\item \textbf{Ethics approval and consent to participate:} [PLACEHOLDER]
\item \textbf{Consent for publication:} Not applicable.
\item \textbf{Author contributions:} [PLACEHOLDER --- to be assigned per Phase C/D progress.]
\end{itemize}

%% References inline (Nature Portfolio convention)
\bibliography{references}

\clearpage

% ===========================================================================
% LIST OF DISPLAY ITEMS  (placeholder — to populate as figures/tables produced)
% ===========================================================================
\section*{List of Display Items}

\subsection*{Main-text Figures}
\begin{description}
\setlength{\itemsep}{2pt}
\item[Figure 1] Yelmen Fig 1 reproduction on 1KG chr21+22 — $\rho_{\max}$ distributions for the diff-chromosome and same-chromosome designs. Sources: \texttt{paper/epistasis\_v1/figures/fig1\_rho\_max\_distribution.\{png,pdf,json\}}.
\item[Figure 2] Motif-typed vs random $Z$, null-target arm — random-$Z$ mean 0.46 vs motif-$Z$ mean 0.45 on $n=3000$ null target SNPs. Sources: \texttt{fig2\_motif\_z\_comparison.\{png,pdf,json\}}.
\item[Figure 3] Motif-typed vs random $Z$, causal-target arm — 2$\times$2 grid (target type $\times$ $Z$ type) demonstrating biology-typed $Z$ surfaces bias that random $Z$ hides ($\Delta\mu=+1.7\sigma$). Sources: \texttt{fig3\_motif\_z\_causal\_arm.\{png,pdf,json\}}.
\item[Figure 4] [PLACEHOLDER] Conservativeness rescue: REGENIE LMM with vs without GraphGWAS-detected interaction covariates.
\item[Figure 5] [PLACEHOLDER] Yeast real-data scenario $\times$ method grid — BCY1$\times$TPK1 recovery rank across 3 phenotype regimes $\times$ 5 methods. Source: \texttt{results/paper2\_epistasis/yeast\_validation.json}.
\item[Figure 6] [PLACEHOLDER] Higher-order: gene-priority RWR vs variant-priority on yeast §Y.4 + canonical-PPI sensitivity panel. Source: \texttt{results/paper2\_epistasis/y5\_gene\_priority\_rwr\_*.json}.
\item[Figure 7] [PLACEHOLDER] Cross-species replication (Arabidopsis FT$\times$FLC, rice Hd1$\times$Hd3a).
\end{description}

\subsection*{Main-text Tables}
\begin{description}
\setlength{\itemsep}{2pt}
\item[Table 1] [PLACEHOLDER] Method portfolio for graph-native epistasis (M1--M7).
\end{description}

\subsection*{Supplementary Figures} \begin{description}\item [Figure S1--Sn] [PLACEHOLDER]\end{description}
\subsection*{Supplementary Tables} \begin{description}\item [Table S1--Sn] [PLACEHOLDER]\end{description}
\subsection*{Supplementary Notes} \begin{description}\item [Note S1--Sn] [PLACEHOLDER]\end{description}

\clearpage

% ===========================================================================
% SUPPLEMENTARY MATERIALS
% ===========================================================================
% PLACEHOLDER supplementary — Notes S1-Sn to populate alongside main text.

\section*{Supplementary Materials}

\subsection*{Note S1: Reproduction of Yelmen \emph{et~al.}\ Fig~1 on 1KG chr21+22}\label{sup:fig1_reprod}

[PLACEHOLDER --- show our $\rho_{\max}$ distributions reproduce theirs
within tolerance; this is the validation step for the
\texttt{graphgwas.bias} implementation.]

\subsection*{Note S2: Mathematical proofs}\label{sup:proofs}

[PLACEHOLDER --- our derivations beyond Yelmen et~al.\ (e.g., the
suppression theorem for motif-typed $Z$ on causal target SNPs);
recap of Theorem~1 (LPCE search-space reduction) from
paper~\#1 \cite{estaji2026graphgwas} Sup.\ Note~S1.]

\subsection*{Note S3: Implementation and architecture}\label{sup:impl}

[PLACEHOLDER --- mirror paper~\#1's Sup.\ Note~S2 layout.]

\subsection*{Note S4: Cross-species ground-truth catalogues for epistasis}\label{sup:catalogues}

[PLACEHOLDER --- yeast literature pairs; human two-variant
heritability catalog; \emph{Arabidopsis} two-locus QTL studies.]


\end{document}
